\chapter{Synthesis and Recommendations}
\label{ch:supp:synthesis}

\section{Measurement precedes allocation}
\label{sec:supp:syn:thread}

The two questions appear distinct but reduce to one structural point. New South Wales is approving large, concentrated, largely inflexible load and contemplating efficiency conditions before it has the measurement and cost-allocation machinery that either task requires. An efficiency-conditioned regime in the Singaporean manner cannot allocate capacity against a Power Usage Effectiveness standard that is neither measured nor disclosed, and a consumer-protection posture cannot prevent energy gentrification without a binding rule on who funds augmentation. Both therefore turn on the same precondition, which the governance literature identifies as the reason compute is a tractable object of policy at all, namely that it is \enquote{detectable, excludable, and quantifiable} \autocite{sastryComputingPowerGovernance2024}. A recurring failure in comparable jurisdictions has been the converse, observed of Nordic datacentre development, where \enquote{no one is steering this development, or closely looking at the impacts that it may have on remote communities in the Arctic and Nordic region} \autocite{sovacoolWholeSystemsEnergy2022}. The remedy in both cases is to measure and disclose first, then allocate and condition.

\section{Consolidated recommendations}
\label{sec:supp:syn:recs}

\begin{enumerate}
\item Make disclosure a condition of consent and continued operation for data centres above a defined information technology power threshold calibrated on the European Union's 500 kilowatt model, covering a measured basket of the NABERS Whole Facility rating, Power Usage Effectiveness, water usage effectiveness, absolute annual energy and renewable share, reported to the regulator and published \autocite{ecEnergyPerformanceDataCentres2024, imanHiddenCostsIntelligence2025}.
\item Convert the Government's stated principle that data centres \enquote{fund their infrastructure requirements} into a binding cost-recovery rule, so that network and water augmentation driven by data centre connection is borne by the connecting party rather than socialised \autocite{inswDataCentreConsultation2026}.
\item Require demonstrably additional renewable generation and firming as a condition of consent, since this is the variable that separates a 26 per cent wholesale price impact from a 3 per cent one and prevents new load leaning on supply built for other consumers \autocite{cefcBaringaDataCentres2025, climateCouncilNSWDataCentres2026, greenpeaceEnergyVampires2026}.
\item Coordinate approvals, transmission planning and renewable energy zone development, which are presently \enquote{undertaken through separate processes across federal and state jurisdictions} and consequently fall out of step \autocite{usscPoweringCloud2025}.
\item Commission and maintain a standing dataset on consumer bill impacts and fleet efficiency, with a Western Sydney regional breakdown, so that policy is set against local measurement rather than overseas analogues \autocite{ecaConsumerImpactsDataCentres2025}.
\item Require disclosure and limitation of backup and off-grid fossil generation, capturing the water and emissions dimensions that a Power Usage Effectiveness figure omits and that have historically gone unmeasured \autocite{climateCouncilNSWDataCentres2026, liMakingAILess2025}.
\end{enumerate}

These measures are mutually reinforcing. Disclosure makes the cost-recovery and renewable-matching rules auditable, the coordination measure prevents the timing failures that raise costs, and the standing dataset closes the evidentiary gap that currently prevents New South Wales from answering the Committee's second question with local numbers. The same instruments also create the option value of treating data centres as flexible grid assets rather than passive loads, an opportunity the literature documents where waste heat is valorised and \enquote{data centre operators in Stockholm and Paris direct waste heat through metropolitan district heating systems and urban homes} \autocite{velkovaDataThatWarms2016}.

\section{Limitations}
\label{sec:supp:syn:limits}

Three limitations should temper how this submission is read. First, several New South Wales figures derive from forecasts and from grey literature rather than from audited operational data, which is itself the strongest argument for the disclosure regime recommended above, since \enquote{the future impact of the data centers' electricity consumption is vulnerable to behavioral usage trends} and is therefore intrinsically uncertain \autocite{kootUsageImpactData2021}. Second, Power Usage Effectiveness is a partial metric, and the water dimension in particular remains poorly measured even internationally, where \enquote{nearly half of servers are fully or partially powered by power plants located within water stressed regions} \autocite{siddikEnvironmentalFootprintData2021a}. Third, the energy gentrification conclusion is bounded by design, asserting a material and preventable risk rather than a demonstrated New South Wales effect, and it should be revisited when the consumer-impact research reports. The inquiry's own final report is the appropriate forum to consolidate that evidence, and the recommendations above are framed so that they can be adopted incrementally as it does.
